(1) 正则概率 \(p_j=e^{-\beta E_j}/Z\)，而 \(\partial_\beta Z=-\sum_j E_je^{-\beta E_j}\)，故平均能量
\[
U=\langle E\rangle=-\partial_\beta\ln Z.
\]
这一恒等式仅需能级与 β 无显式依赖，且配分函数收敛。

(2) 再对 β 求导，
\[
\partial_\beta^2\ln Z
=\langle E^2\rangle-\langle E\rangle^2
=\mathrm{Var}(E).
\]
由于 \(\partial_\beta=-k_BT^2\partial_T\)，有
\[
\mathrm{Var}(E)=k_BT^2\left(\frac{\partial U}{\partial T}\right)_V
=k_BT^2C_V\ge0.
\]
若假定正则平衡且能级独立于 T，则热容不能为负。

(3) 对 N 个弱相关单元，\(U=O(N)\)、\(C_V=O(N)\)，所以标准差 \(\Delta E=\sqrt{k_BT^2C_V}=O(\sqrt N)\)，相对涨落 \(\Delta E/U=O(N^{-1/2})\)。这解释宏观体系虽然能量在正则系综中不严格固定，实际相对变化却极小。

(4) 当 T→0 并具有非简并基态且有限激发能隙时，能量分布集中在基态，方差趋零；公式中的 \(T^2C_V\to0\)。对于相变或长程相互作用等不满足普通热力学极限的体系，\(C_V\) 可能超线性增长，必须重新检查相对涨落的标度。
